% !TEX program = pdflatex
\documentclass[11pt]{article}
\usepackage[T1]{fontenc}
\usepackage{lmodern,amsmath,amssymb,amsthm,microtype,booktabs}
\usepackage[margin=1in]{geometry}
\usepackage[hidelinks]{hyperref}
\hypersetup{pdftitle={A cubic higher multiplication creates infinitely many bounded profiles}}
\newtheorem{thm}{Theorem}
\newtheorem{lem}[thm]{Lemma}
\newtheorem{cor}[thm]{Corollary}
\newcommand{\Dfd}{D_{\rm fd}}
\newcommand{\per}{\operatorname{per}}
\newcommand{\UU}{\mathcal U}
\newcommand{\thick}{\operatorname{thick}}
\newcommand{\Ext}{\operatorname{Ext}}
\newcommand{\Hom}{\operatorname{Hom}}
\newcommand{\Cone}{\operatorname{Cone}}
\newcommand{\RHom}{\operatorname{RHom}}
\setlength{\emergencystretch}{1em}
\title{A cubic higher multiplication creates\\ infinitely many bounded profiles}
\author{Marouane Lamharzi Alaoui\\
\small Research notes on two-term dg kernels and bounded profiles of universal localisations}
\date{18 November 2026}
\begin{document}
\maketitle
\begin{abstract}
We exhibit a six-dimensional commutative local differential graded algebra
over any field. Its formal cohomology algebra has two bounded profiles.
The nonformal algebra admits an infinite antichain, including over finite
fields. Six is the least possible total dimension for a change of poset
among strict two-term models with local degree-zero cohomology and residue
field $k$. More precisely, reverse inclusion of finite sets of closed points
of the affine line embeds in its profile poset. A second, period-two
detector gives an additional intermediate profile outside that embedding.
Combined with the manuscript's realisation theorem, this changes the
marked bounded defects of ordinary universal localisations at fixed
graded kernel cohomology.
\end{abstract}

\section{The pair}
Use left dg modules and cochain grading over a field $k$. For
$E\in\Dfd(D)$ set
\[
 \UU_E=\{X\in\Dfd(D)\mid \dim_k\bigoplus_n\Ext_D^n(X,E)<\infty\}.
\]
The marked bounded defect is $\UU_E/\per D$. Each profile is thick and
contains $\per D$, so passage to this quotient preserves and reflects
inclusion. Write $\mathcal P$ for the monic irreducibles in $k[T]$,
equivalently the closed points of $\mathbb A^1_k$.

\begin{thm}\label{main}
There is a commutative two-term local dg algebra $D$ of total dimension
six such that $F=H^*(D)$ with zero differential has exactly two
profiles, while the profile poset of $D$ contains an order-embedded copy
of $(\{L\subset\mathcal P:\,L\text{ finite}\},\supseteq)$.
There is also a proper intermediate profile outside this copy.
Both singularity categories are nonzero.
\end{thm}

Take $D^0=\langle1,a,b,s\rangle$, $D^{-1}=\langle r,w\rangle$.
The complete list of nonzero radical products and differentials is
\[
             a^2=s,\qquad br=rb=w,\qquad dr=s .
\]
The radical cube vanishes, $s$ kills the radical, and negative elements
multiply to zero. This checks associativity, graded commutativity, and
the Leibniz rule in every characteristic. Thus $H^0(D)=k\oplus ka\oplus kb$,
$H^{-1}(D)=kw$, with all radical products in cohomology zero.
In homotopy transfer take $h(s)=r$ and $h$ zero on the other basis
elements. In the convention $m_3(i,j,k)=i h(jk)-h(ij)k$ on degree-zero
radical inputs, the two nonzero values are
\[
              m_3(b,a,a)=w,\qquad m_3(a,a,b)=-w .
\]
Both survive in characteristic two. Every normalized Hochschild coboundary
in the relevant degree vanishes on reduced triples, because the radical
squares to zero and kills $kw$. Higher arities vanish by degree.
Consequently a single nonzero ternary operation, with the two coefficients
required by this model, gives the difference from $F$.

This is a jump in a flat one-parameter family. Over $k[t]$ retain the
same six free basis vectors and differential, and set $br=rb=t w$.
The underlying complex splits into its free cohomology plus the
contractible pair $r\mapsto s$, and its graded cohomology algebra is
$F\otimes_k k[t]$. At $t=0$ the projection killing $r,s$ is a dg
quasi-isomorphism to $F$. In any field-valued fibre with $t\ne0$,
rescaling $w$ identifies the fibre with the base change of $D$.

\section{The Koszul dual and its normal words}
Let $x,z$ be dual to $a,b$, in degree one. The dual of the shifted
class $w$ has degree two and is denoted $y$. The reduced bar construction
on the augmented minimal model gives a model for $\RHom_D(k,k)$,
\[
 K=k\langle x,z,y\rangle,\qquad dx=dz=0,\qquad dy=x^2z-zx^2 .
\]
Indeed dualising the two cubic coefficients gives the displayed
commutator up to an overall sign, absorbed into $y$. Reversing the
word convention changes it by a sign as well. The grading is positive
on the generators, so degreewise duality introduces no completed
infinite sums in any fixed degree.

One can check the sign from the strict model before transfer. Write $u,v$
for the shifted duals of $s,r$, in degrees one and two. With rescaled
generator signs the strict reduced bar dual has
$du=v-x^2$, $dv=0$, and $dy=vz-zv$, with $dx=dz=0$. For example, in
the suspended bar convention $b_1([i])=-[di]$ and
$b_2([i|j])=(-1)^{|i|}[ij]$, the terms $[b|r]$ and $[r|b]$ give
opposite signs. Replacing $v$ by the generator $v-x^2$ and $y$ by
$y-(uz+zu)$ splits off a free pair with $du=v-x^2$; the new $y$ has
differential $x^2z-zx^2$. The kernel of the projection deleting the pair
is acyclic: split it into words alternating the original generators'
subalgebra with positive tensor powers of the contractible pair.
Each summand contracts over $k$.

\begin{lem}\label{normal}
The map $K\to A=k\langle x,z\rangle/(x^2z-zx^2)$ sending $y$ to zero
is a quasi-isomorphism for $A$ with zero differential. The words in
$x,z$ avoiding $xxz$ form a basis of $A$. In particular its series is
$1/((1-t)(1-t-t^2))$.
\end{lem}
\begin{proof}
Give $x,z$ weight one and $y$ weight three. At each fixed weight filter
words by lexicographic order of the word obtained by expanding $y$ to
$xxz$, with $x>z$. Replacing $y$ by $zxx$ strictly lowers that order.
The associated graded differential replaces $y$ by $xxz$ with its
cochain sign. Fix an expanded word. No two occurrences of $xxz$ overlap,
since it has no proper nonempty suffix equal to a prefix. Each occurrence
can independently appear as $y$ or as $xxz$. Its two choices give a
two-term complex with differential an isomorphism, with tensor product
signs from the degrees of the fixed intervening letters and the other
choices. The subcomplex contracts if there is any occurrence. Otherwise
one basis word survives in degree equal to the weight. The finite
filtration therefore has cohomology concentrated in that degree.
All surviving classes lift to the indicated words in $x,z$. In $A$ the
rewrite $xxz\mapsto zxx$ terminates in fixed length. Its only possible
simultaneous rewrites are disjoint and commute, proving uniqueness of
normal form. The projection identifies these two bases. Every normal word
is a sequence of blocks $z,xz$ followed by a power of $x$, giving the series.
\end{proof}

The functor $\RHom_D(k,-)$ identifies $\Dfd(D)=\thick(k)$ with the
perfect \emph{right} dg modules over $A$. Indeed the functor takes the
generator to the free rank-one module and induces the isomorphism on its
graded endomorphisms. Full faithfulness and essential surjectivity on
the thick closures follow by cones and summands. Here $\thick(k)=\Dfd(D)$
by truncation and finite length in the heart. In particular the finite
free cones constructed below correspond to actual finite-cohomology
$D$-modules even though their cohomology on the $A$ side is infinite.
The use of right modules makes matrix entries act on free columns by
left multiplication and on detector coefficients by right multiplication.

\section{Detectors and witnesses}
First fix $\lambda\in k$ and set $u=z-\lambda x$. The relation becomes
$x^2u=ux^2$. The normal words from Lemma~\ref{normal} work in $x,u$
as well. Give $M_\lambda=k[x]$ the graded right action
$z\mapsto\lambda x$. There is an exact sequence of graded right modules
\[
 0\longrightarrow A[-1]\oplus A[-2]
  \xrightarrow{(u,\,xu)} A\longrightarrow M_\lambda\longrightarrow0 .
\]
A normal word involving $u$ begins either with $u$ or with $xu$.
Prefixing either block to any normal word produces a normal word, and
the two images are disjoint. This proves exactness and makes $M_\lambda$
perfect as a dg $A$-module. A shift $[-j]$ places a generator in degree $j$.

For $p\in\mathcal P$, write $k_p=k[T]/(p)$, $\alpha_p=T\bmod p$,
$d_p=\deg p$. Use the same construction over $k_p$ at $\lambda=\alpha_p$,
then restrict scalars to $A$. Denote the resulting right module by
$M_p=k_p[x]$. It is perfect, because $A\otimes_k k_p$ is finite free over
$A$. It has dimension $d_p$ in every nonnegative degree. Define
$E_p\in\Dfd(D)$ to be its preimage under the equivalence above.
For a monic irreducible $q$ let $T_q$ be the matrix of multiplication by
$T$ on $k_q$ in its power basis, and let $X_q$ correspond to
\[
 C_q=\Cone\big(A[-1]^{d_q}
         \xrightarrow{\ z I-xT_q\ } A^{d_q}\big).
\]

\begin{lem}\label{incidence}
$X_q\in\UU_{E_p}$ if and only if $q\ne p$.
\end{lem}
\begin{proof}
Apply $\RHom_A(-,M_p)$ to the defining triangle of $C_q$. In sufficiently
high degrees the induced map on the two free terms is multiplication by
$x$ times the constant matrix $\alpha_p I-T_q^{\mathsf t}$ over $k_p$,
up to a common cochain sign and the choice of row convention. Multiplication
by $x$ identifies successive homogeneous components of $M_p$. The matrix
has determinant $q(\alpha_p)$, zero precisely for $q=p$. In that case
the kernel and cokernel persist in every high degree, and otherwise both
vanish there. The low degrees are finite and the complexes are bounded
below. The test works for inseparable residue extensions as well.
\end{proof}

For finite $L\subset\mathcal P$ take $E_L=\bigoplus_{p\in L}E_p$,
including $E_{\varnothing}=0$. Its profile is the intersection of the
individual profiles. By Lemma~\ref{incidence}, $X_q\in\UU_{E_L}$ exactly
when $q\notin L$. It follows that $\UU_{E_L}\subseteq\UU_{E_{L'}}$
exactly when $L\supseteq L'$. There are infinitely many monic irreducibles
over every field, by the product-plus-one argument. The singleton sets
give the infinite antichain. Each nonempty $L$ gives a proper profile,
because $\Ext_D^*(k,E_L)$ is infinite. It also contains a nonperfect
object: take $q\notin L$; $X_q$ is nonperfect since it has infinite Ext
to the finite detector $E_q$.

As a check on the role of the witnesses, replacing $X_q$ by $E_q$ in
Lemma~\ref{incidence} would give a false statement. For scalar points
$\lambda,\mu$, the two-term free resolution of $M_\lambda$ gives in
every high degree a map from one copy of $k$ to two copies of $k$ after
applying $\RHom_A(-,M_\mu)$. Its cokernel is always nonzero. In particular
$\Ext_D^*(E_\lambda,E_\mu)$ is unbounded even at distinct scalar points.

For $F$, the augmentation ideal $J=ka\oplus kb\oplus kw$ has $J^2=0$.
Lemma 1 of the companion note shows that every finite object modulo
perfects is a sum of shifted copies of $k$. Briefly, cut a bounded-above
minimal resolution below all cohomology, discarding a finite cell
submodule $P'$. In the quotient, what remains is a shift of the smart
truncation at the cut of $P'$. Both terms lie in $JP'$, so this is just
a complex of vector spaces. Every nonzero object therefore generates the
quotient as a thick subcategory. The zero and residue-field detectors
give both endpoints, and $k$ survives because its formal reduced bar
dual is the free algebra on generators of degrees $1,1,2$. This proves
the two-profile assertion.

\section{A detector outside the closed-point family}
Let $N$ have one-dimensional graded pieces $k e_n$, $n\ge0$,
and actions $e_n x=e_{n+1}$, $e_n z=\lambda_n e_{n+1}$, where
$\lambda_{2i}=0$, $\lambda_{2i+1}=1$. It is a right $A$-module since
$\lambda_{n+2}=\lambda_n$. We claim it has the resolution
\[
 0\longrightarrow A[-1]\oplus A[-2]
 \xrightarrow{(z,\,xz-x^2)}A\xrightarrow{1\mapsto e_0}N\longrightarrow0.
\]
The indicated two relations annihilate $e_0$. For injectivity, take
nonzero inputs and filter normal words by the number of $z$'s. In maximal
input filtration degree the leading outputs are $z$ and $xz$ prefixed to
normal words, disjoint sets. The correction $-x^2$ cannot increase the
number of $z$'s under normalisation. Thus the two-column map is injective.
By Lemma~\ref{normal} its graded cokernel has series
$(1-t-t^2)H_A(t)=1/(1-t)$; the surjection onto $N$ is an isomorphism.
Let $E_N$ be the object's preimage. Test it by the cone $C$ of left
multiplication by $zx-xz$ from $A[-2]$ to $A$. On $M_p$ this commutator
acts by zero for every $p$. On $N$ it acts in degree $n$ by
$(\lambda_n-\lambda_{n+1})e_{n+2}$, always an isomorphism, including
in characteristic two. The preimage of $C$ lies in $\UU_{E_N}$ and
outside every $\UU_{E_L}$ for nonempty $L$; it is nonperfect since
its Ext to $E_p$ is unbounded. Also $k\notin\UU_{E_N}$, since $N$ has
infinite cohomology. This proves the last assertion of Theorem~\ref{main}.

\section{Sharpness of the dimension}
\begin{cor}
Let $D'$ be a finite-dimensional unital dg algebra supported in degrees
$-1,0$ over $k$, with $H^0(D')$ local and residue field $k$. If
$\dim_k D'\le5$, its bounded-profile poset is isomorphic to the poset
of its formal cohomology algebra. Thus six is minimal for the change in
Theorem~\ref{main}, even among associative models in this class.
\end{cor}
\begin{proof}
Write $h_i=\dim H^i(D')$ and $s=\operatorname{rank}d$. Then
$\dim D'=h_0+h_{-1}+2s$. With $s=0$ the differential is zero, and
$h_{-1}=0$ likewise implies formality by projection to $H^0$. If $h_0=1$,
a strictly unital minimal model has no reduced degree-zero inputs; degree
considerations exclude all higher operations. The only remaining case in
dimension at most five has $(h_0,h_{-1},s)=(2,1,1)$. The degree-zero
algebra is the dual numbers and its radical kills the one-dimensional
negative cohomology on both sides. The sole possible higher coefficient
is $m_3(t,t,t)=\gamma w$. For $\gamma=0$ it is formal, and all nonzero
values are equivalent by rescaling $w$. Thus the nonzero case has the
augmented derived type of Example 3.11. Theorem 3 of the companion note
classifies both posets there as two-element chains. This proves the bound.
\end{proof}

\section{Realisation, scope and prior work}
The matrix construction of Theorem 3.7 of the manuscript uses $\dim
R=3\dim D^0+\dim D^{-1}$ and target dimension $4\dim D^0$. It gives
a full-profile ordinary localisation of dimensions $14\to16$ for $D$
and $10\to12$ for $F$. Corollary 3.10 then realises each detector
profile as a marked bounded defect for its respective kernel. The
dimension counts refer to the initial full-profile maps.

Geometric detection of pairwise Ext vanishing is classical for ordinary
complete intersections, as in \href{https://doi.org/10.1007/s002220000090}{Avramov and Buchweitz}, and
\href{https://arxiv.org/abs/1912.12009}{Pollitz} treats pairs of dg modules
over Koszul complexes. In particular, his Proposition 3.3.1 gives finite
generation over degree-two polynomial operators for a derived complete
intersection. Applied to its residue field this bounds Betti numbers by
a polynomial. Lemma~\ref{normal} instead gives exponential growth, which
excludes that derived-complete-intersection case for our augmented model.
The point distinction in Lemma~\ref{incidence}
has that character, although it is proved directly from finite free
resolutions over the displayed one-relator algebra. The Koszul-dual
calculation uses the classical nonoverlap mechanism in
\href{https://doi.org/10.1016/0021-8693(82)90104-1}{Anick's work};
the link between higher products and relation coefficients has precedents
in \href{https://arxiv.org/abs/math/0606144}{Lu, Palmieri, Wu and Zhang}.
The formal endpoint is consistent with the rigidity of split
commutative dg extensions studied by
\href{https://arxiv.org/abs/1508.00748}{Avramov, Iyengar, Nasseh and
Sather-Wagstaff}. The comparison at a fixed graded cohomology algebra,
the explicit six-dimensional strict model, and the passage to marked
defects are the claims for which priority still requires review.

\paragraph{Starting point for a classification.}
The element $x^2$ is central. There is an exact presentation of its
graded localisation,
\[
 A[(x^2)^{-1}]\cong k\langle U,V\rangle[X,X^{-1};\sigma],
 \qquad \sigma(U)=V,\quad\sigma(V)=U .
\]
The skew relation is $Xf=\sigma(f)X$. Send $x$ to $X$ and $z$ to
$UX$. Conversely take $U=zx^{-1}$ and $V=xzx^{-2}$; conjugation by
$x$ interchanges them because $x^2$ is central. These formulas are
mutually inverse. In particular the degree-zero algebra in this
localisation is the free associative algebra on two generators.

Every pair of operators $T_0,T_1$ on a finite-dimensional space $W$
already gives a perfect detector on the original $A$ side. Put $W$ in
each nonnegative degree, let $x$ act by the identity between successive
copies and let $z$ act by $T_0,T_1$ alternately. The degree-zero
generators have degree-one relations $v z-(T_0v)x$ and degree-two
relations $v xz-(T_1v)x^2$ for $v$ in a basis of $W$. The filtered
injectivity proof for $N$ works verbatim: the two leading prefixes are
disjoint and matrix coefficients in the correction terms introduce no
$z$. The cokernel has series $\dim(W)/(1-t)$ and surjects onto the
specified module. This proves perfection for all such pairs.
Determining when distinct pairs give the same profile, and accounting
for perfect complexes whose cohomology does not have eventually
invertible $x$, remain open here. The explicit extra profile in Section 4
shows that the finite closed-point family alone is insufficient.
\end{document}
